% ===========================================================================
%  Annexe réservée au professeur — progression annuelle, sans dates
%  Incluse uniquement dans l'édition professeur (voir livre-t12ose.tex).
% ===========================================================================
\chapter{Progression annuelle}
\markboth{Progression annuelle}{Progression annuelle}

Cette annexe ne porte aucune date, pour la même raison que dans le livre de
série S : le calendrier scolaire change chaque année, une période est
parfois amputée par les examens blancs ou par les intempéries, mais l'ordre
des chapitres et les volumes horaires, eux, ne bougent pas. Reportez vos
propres dates dans la dernière colonne au début de l'année.

La série OSE pose une difficulté que la série S ne pose pas : le Programme
d'Études 2026 décrit un bloc de géométrie plane complet — nombres complexes,
barycentre, transformations, similitudes —, mais la Répartition Annuelle du
Programme d'Études 2025-2026, celle qui fixe le calendrier réellement suivi
en classe, ne le mentionne à aucun moment pour la Terminale OSE. La
progression proposée ici tranche cet écart en insérant le bloc de géométrie
là où il perturbe le moins le reste de l'année ; la section « Le bloc
Géométrie » qui suit l'explique en détail. Un professeur dont l'établissement
suit une Répartition Annuelle qui ignore encore ce bloc peut le traiter en
fin d'année, ou l'alléger, sans remettre en cause le reste du calendrier.

\section*{Vue d'ensemble}

L'année compte cinq périodes. Le tableau ci-dessous fait la conversion entre
l'ordre d'enseignement proposé et le rangement du livre, qui suit les
domaines.

\vspace{0.6em}
\begin{center}
\setlength{\tabcolsep}{6pt}
\renewcommand{\arraystretch}{1.5}
\begin{tabular}{@{}L{1.3cm} L{6.5cm} L{1.6cm} L{2.5cm}@{}}
\toprule
\textsf{\footnotesize\bfseries PÉRIODE} &
\textsf{\footnotesize\bfseries CHAPITRES} &
\textsf{\footnotesize\bfseries HEURES} &
\textsf{\footnotesize\bfseries VOS DATES} \\
\midrule
\textsf{\bfseries I}   & 8 Systèmes linéaires \newline 1 Limites et continuité \newline 2 Dérivabilité \newline 3 Étude et représentation graphique & 48 h & \\
\textsf{\bfseries II}  & 4 Primitives et calcul intégral \newline 5 Fonctions logarithmes \newline 9 Nombres complexes & 43 h & \\
\textsf{\bfseries III} & 6 Exponentielle et fonctions puissances \newline 7 Suites numériques \newline 10 Barycentre & 37 h & \\
\textsf{\bfseries IV}  & 13 Probabilités \newline 14 Variables aléatoires et lois \newline \textit{Annexe : Mathématiques financières} & 42 h* & \\
\textsf{\bfseries V}   & 15 Statistique à deux variables \newline 11 Transformations et isométries \newline 12 Similitudes planes & 30 h & \\
\midrule
& \textsf{\bfseries Total} & \textsf{\bfseries 200 h} & \\
\bottomrule
\end{tabular}
\end{center}

\vspace{0.4em}
\noindent Cent quatre-vingts heures correspondent exactement aux quatre
composantes chiffrées du Programme d'Études — quatre-vingts pour l'analyse,
vingt-cinq pour l'algèbre, quarante-cinq pour la géométrie, trente pour les
données et probabilités : ce sont elles qui forment les quinze chapitres
numérotés du livre. Les vingt heures marquées d'un astérisque en période IV
couvrent l'annexe de mathématiques financières, absente du chiffrage
officiel du Programme d'Études mais bien réelle : la Répartition Annuelle
l'installe précisément en période 4, et l'épreuve du baccalauréat lui
consacre chaque année un exercice entier. C'est pour cette raison — hors du
chiffrage officiel — que ce livre la présente en annexe plutôt que comme une
cinquième partie numérotée, sans pour autant en réduire le contenu ni le
temps qui doit lui être consacré.

\section*{Découpage indicatif en séances}

Le découpage suppose des séances de deux heures. Quatre chapitres — 7, 9 et,
dans une moindre mesure, d'autres — laissent une heure qui ne se divise pas
exactement ; absorbez-la dans la séance de révision qui suit, ou dans un
devoir court.

\vspace{0.5em}
\begin{center}
\setlength{\tabcolsep}{6pt}
\renewcommand{\arraystretch}{1.35}
\begin{tabular}{@{}L{0.7cm} L{5.3cm} L{1.5cm} L{1.7cm} L{2.1cm}@{}}
\toprule
& \textsf{\footnotesize\bfseries CHAPITRE} &
\textsf{\footnotesize\bfseries HEURES} &
\textsf{\footnotesize\bfseries SÉANCES} &
\textsf{\footnotesize\bfseries ÉVALUATION} \\
\midrule
1  & Limites et continuité                     & 14 h & 7 & devoir court \\
2  & Dérivabilité                              & 12 h & 6 & devoir court \\
3  & Étude et représentation graphique         & 10 h & 5 & — \\
4  & Primitives et calcul intégral             & 16 h & 8 & devoir surveillé \\
5  & Fonctions logarithmes                     & 14 h & 7 & — \\
6  & Exponentielle et fonctions puissances     & 14 h & 7 & devoir surveillé \\
7  & Suites numériques                         & 13 h & 6 & devoir court \\
8  & Systèmes linéaires (pivot de Gauss)       & 12 h & 6 & — \\
\midrule
9  & Nombres complexes                         & 13 h & 6 & devoir surveillé \\
10 & Barycentre                                & 10 h & 5 & — \\
11 & Transformations et isométries             & 12 h & 6 & devoir court \\
12 & Similitudes planes                        & 10 h & 5 & — \\
\midrule
13 & Probabilités                              & 12 h & 6 & devoir court \\
14 & Variables aléatoires et lois              & 10 h & 5 & devoir surveillé \\
15 & Statistique à deux variables               & 8 h  & 4 & — \\
\bottomrule
\end{tabular}
\end{center}

\subsection*{Et l'annexe de mathématiques financières ?}

Hors chiffrage officiel, donc hors du tableau ci-dessus, mais bien présente
au calendrier réel : la Répartition Annuelle la place en période IV, aux
côtés des probabilités. Le découpage indicatif est le même que pour un
chapitre ordinaire.

\vspace{0.5em}
\begin{center}
\setlength{\tabcolsep}{6pt}
\renewcommand{\arraystretch}{1.35}
\begin{tabular}{@{}L{6.0cm} L{1.5cm} L{1.7cm} L{2.1cm}@{}}
\toprule
\textsf{\footnotesize\bfseries VOLET} &
\textsf{\footnotesize\bfseries HEURES} &
\textsf{\footnotesize\bfseries SÉANCES} &
\textsf{\footnotesize\bfseries ÉVALUATION} \\
\midrule
Intérêts et escompte    & 10 h & 5 & devoir court \\
Annuités et emprunts    & 10 h & 5 & devoir surveillé \\
\bottomrule
\end{tabular}
\end{center}

\section*{Le bloc Géométrie et la Répartition Annuelle}

Un mot d'explication mérite d'être écrit noir sur blanc, plutôt que laissé à
la découverte. Le Programme d'Études 2026 consacre quarante-cinq heures à un
bloc de géométrie plane — nombres complexes, barycentre, transformations et
isométries, similitudes — presque aussi riche que celui de la série S. La
Répartition Annuelle 2025-2026, qui fixe le calendrier réellement suivi dans
les établissements, ne le mentionne nulle part pour la Terminale OSE : ses
cinq périodes ne couvrent que l'algèbre, l'analyse, les suites, les
mathématiques financières, la probabilité et la statistique à deux
variables. Le sujet de baccalauréat de la session 2026, seul disponible en
intégralité au moment de la rédaction de ce livre, confirme cet écart : ni
nombres complexes, ni barycentre, ni similitudes n'y apparaissent.

Ce livre choisit de traiter le bloc entier, pour deux raisons. La première
est que le Programme d'Études reste le texte de référence, et qu'une session
future peut très bien rattraper ce qu'elle n'a pas encore interrogé — un
professeur qui aurait sauté ce bloc se trouverait démuni le jour où cela
arrive. La seconde est que ces notions forment un ensemble cohérent et
transférable, utile bien au-delà de l'épreuve elle-même. Le tableau plus
haut le glisse en trois morceaux — nombres complexes en période II,
barycentre en période III, transformations et similitudes en période V —
choisis pour ne jamais rompre l'ordre interne du bloc lui-même : les
nombres complexes sont posés avant qu'on écrive une transformation sous
forme complexe, et les transformations avant les similitudes qui les
généralisent.

Un établissement dont la Répartition Annuelle locale ignore encore ce bloc
peut légitimement l'alléger — se limiter aux nombres complexes et au
barycentre, par exemple, et renvoyer transformations et similitudes à une
révision de fin d'année — sans mettre en péril la préparation à l'épreuve
telle qu'elle se pose aujourd'hui.

\section*{Les repères de fin de période}

Ce que l'élève doit savoir faire avant de passer à la suite.

\subsection*{À la fin de la première période}

Résoudre un système de trois équations à trois inconnues par la méthode du
pivot de Gauss, en rédigeant chaque étape d'élimination. Déterminer une
limite, y compris celle d'une composée, et rédiger complètement une
application du théorème des valeurs intermédiaires. Étudier la dérivabilité
en un point et reconnaître un point anguleux. Mener une étude de fonction
jusqu'au tracé, asymptotes et position relative comprises.

\subsection*{À la fin de la deuxième période}

Reconnaître une forme $u'u^{n}$ ou $\frac{u'}{u}$ et calculer une intégrale
par parties. Résoudre équations et inéquations faisant intervenir le
logarithme népérien, en posant les conditions d'existence avant tout calcul.
Passer d'une forme algébrique à une forme trigonométrique d'un nombre
complexe et résoudre une équation du second degré dans $\C$.

\subsection*{À la fin de la troisième période}

Reconnaître les limites usuelles de l'exponentielle et des fonctions
puissances, en particulier les croissances comparées. Démontrer une
propriété par récurrence et établir la convergence d'une suite. Construire
un barycentre et déterminer une ligne de niveau.

\subsection*{À la fin de la quatrième période}

Calculer un intérêt simple, un escompte commercial, une valeur acquise à
intérêt composé, et passer de l'un à l'autre sans confondre les taux.
Construire un tableau d'amortissement et retrouver l'annuité constante d'un
emprunt. Construire un arbre de probabilité, calculer une probabilité
conditionnelle, donner la loi d'une variable aléatoire, reconnaître une
situation binomiale.

\subsection*{À la fin de la cinquième période}

Calculer un coefficient de corrélation et une droite de régression par la
méthode des moindres carrés, et l'utiliser pour une estimation. Décomposer
une isométrie en réflexions et identifier une similitude à partir de son
expression complexe.

\section*{Les points de vigilance}

Quatre passages méritent une attention particulière.

Le premier est méthodologique : la Répartition Annuelle mentionne la
méthode de Cramer pour le système $3\times3$, tandis que le Programme
d'Études — et donc ce livre — enseigne le pivot de Gauss. Les deux méthodes
sont également valables et conduisent au même résultat ; il est utile de
signaler brièvement l'existence de Cramer aux élèves qui la rencontreraient
ailleurs, sans en faire un second cours.

Le deuxième concerne les mathématiques financières : c'est, pour la plupart
des élèves, un vocabulaire entièrement nouveau — valeur nominale, valeur
actuelle, escompte, agios, amortissement. Une séance d'introduction au seul
vocabulaire, avant la première formule, évite bien des confusions
ultérieures entre intérêt et escompte, ou entre annuités constantes et
amortissements constants.

Le troisième est le bloc géométrique lui-même, détaillé plus haut : son
insertion demande d'être annoncée aux élèves comme un choix du livre, pas
comme une évidence du programme suivi par leur établissement.

Le quatrième, enfin, est la cinquième période, qui porte la charge la plus
légère en heures mais tombe au moment où la pression des révisions
générales est la plus forte. C'est un bon endroit pour placer, si le temps
manque ailleurs, les séances de synthèse consacrées au recueil de sujets.
